\section{Εικόνες και Πίνακες}
Εικονογραφήσεις μπορούν να σας βοηθήσουν να αποσαφηνιστούν δύσκολες έννοιες. Οι πίνακες χρησιμοποιούνται για την δομημένη απεικόνιση των δεδομένων. Και τα δύο αποτελούν σημαντικά στοιχεία της γραπτής έκθεσης.

\begin{figure}[htb]
\centering
 \includegraphics[width=70mm]{images/teipirlogo.eps}
\caption{Παράδειγμα εισαγωγής μιας εικόνας με λεζάντα}
\label{figure:teipirlogo}
\end{figure}

Ενσωματώστε τις εικόνες και τους πίνακες μες το κείμενο. Τοποθετήστε τις όσο το δυνατόν πλησιέστερα στο σημείο όπου αναφέρθηκαν για πρώτη φορά. Και τα δύο οι απεικονίσεις και οι πίνακες πρέπει να διαθέτουν μια λεζάντα που την βάζουμε κάτω από το τον πίνακα ή εικόνα. Και για τα δύο στοιχεία είναι μια συνεχής, ξεχωριστή αρίθμηση απαραίτητη. Αν στο κείμενο αναφέρεστε σε μια εικόνα ή έναν πίνακα, μπορείτε να το κάνετε μέσω του αριθμού τους. Ένα παράδειγμα γι'αυτό μπορείτε να δείτε σε αυτό το έγγραφο στο τμήμα \ref{sec:έκταση}

\section{Οι παραπομπές}
Σε μεγάλα έγγραφα θα πρέπει να χρησιμοποιείτε πιο έντονα την αρίθμηση των τμημάτων. Στις παραπομπές θα αναφέρετε τον αριθμό του τμήματος. Παράδειγμα: 

\begin{center}
\line(1,0){350}\\
Όπως φαίνεται στο τμήμα \ref{sec:έκταση},  ... \\
\line(1,0){350}\\
\end{center}

Αν δεν χρησιμοποιείτε αριθμούς ενότητας, μπορείτε να αναφερθείτε στον τίτλο της ενότητας. Παράδειγμα: 

\begin{center}
\line(1,0){350}\\
Όπως φαίνεται στο κεφάλαιο ``Εννοιολογικά θεμέλια'', ...\\
\line(1,0){350}\\
\end{center}

\section{Παραπομπές και οι αναφορές}
Σοβαρές επιστημονικές εργασίες αναφέρονται πάντα σε άλλα έργα. Επίσης, συχνά αποσπάσματα κειμένου, οι εικόνες ή πίνακες παρατίθενται αμετάβλητα (αυτούσια). Τέτοιου είδους αποσπάσματα θα πρέπει αυστηρά να αναφέρονται με παραπομπές και τσιτάτα. Η χρήση του ξένου υλικού χωρίς υπόδειξη δεν είναι μόνο μορφολογικά επιβλαβής, αλλά μπορεί σε ορισμένες περιπτώσεις να θεωρηθεί ως μια απόπειρα απάτης.
Είναι σημαντικό να διασφαλιστεί ότι αναφορές να εντάσσονται αρμονικά από τον συντάκτη στο γενικό πλαίσιο της εργασίας του. Αποσπάσματα δεν τα αφήνουμε ποτέ μόνα τους. Θα πρέπει να επεξηγηθούν και να ερμηνευτούν πάντα σε σχέση με την ίδια την εργασία. 
Η αναφορά σε άλλες εργασίες δημιουργείται μέσω  των αναφορών βιβλιογραφίας. 
Υπάρχουν διάφορα ξενόγλωσσα πρότυπα όπως το DIN 1505-2 που διέπει τη σωστή μορφή των βιβλιογραφικών αναφορών στις γερμανόφωνες χώρες. ... Τελευταία η χρήση των ηλεκτρονικών αναφορών πχ. του διαδικτύου θεωρείται επίσης ως έγκυρη βιβλιογραφία για την εργασία. ... Ωστόσο, είναι πολύ σημαντικό το γεγονός ότι η βιβλιογραφία και οι αναφορές στο κείμενο σας θα περιέχουν όλα τα στοιχεία χρειάζονται και θα ακολουθούν κάποιο ένα ενιαίο πρότυπο.
Εάν χρησιμοποιείτε OpenOffice προϊόντα για να δημιουργήσετε το κείμενό σας, δεν θα βρείτε κάποια ικανοποιητική υποστήριξη στην προετοιμασία και χρήση των βιβλιογραφιών. Είστε μόνοι σας. Σε αυτή την περίπτωση, θα πρέπει να κατευθυηθείτε σε κάποι (συνήθως αμερικάνικο) πρότυπο. Εάν χρησιμοποιείτε \TeX{} ή \LaTeX{} ή συναφή συστήματα στοιχειοθεσίας, θα λάβετε μια εκτενή στήριξη για τη δημιουργία βιβλιογραφιών. Το \LaTeX{} και το συνδεδεμένο εργαλείο Bib\TeX{} θα σας δώσει την ευκαιρία να δημιουργήσει μια βιβλιογραφική βάση δεδομένων. Απ' αυτή θα παραχθεί με τη χρήση διαφόρων προτύπων σχετικός κατάλογος αναφορών καθώς και μέσω των  τύπων αναφοράς και αναφορές μέσ' το κείμενο. Είναι σημαντικό να υπάρχουν όλα τα σχετικά δεδομένα σε μια αναφορά βιβλιογραφίας και να χρησιμοποιείται με μια συνεπή μορφή.
Τα αυτούσια τσιτάτα θα πρέπει να μπαίνουν σε εισαγωγικά και να διαχωρίζονται από το τρέχον κείμενο. Σε περίπτωση που δεν χρησιμοποιείται όλο το τσιτάτο αλλά κάοια αποσπάσματα τότε θα πρέπει να παρεμβάλλεται το εξής ``[...]'' ως ένα σημάδι της παράλειψης.

\begin{table}
\begin{tabular}{|l|}\hline
``Υπάρχουν δυο υλοποιήσεις νημάτων: Kernel-Threads και\\
User-Threads. [...] Κάθε νήμα μέσα σε μια διεργασία \\
ξεχωρίζει από έναν μοναδικό αριθμό (Thread ID) που \\
είναι τύπου pthread\_t'', \cite{oslab-threads}\\
\hline
\end{tabular}
\caption{Παράδειγμα 1. Απόσπασμα με αποσιωπητικά}
\end{table}

Όταν χρησιμοποιούνται τμήματα ξένων κειμένων, εικόνων ή πινάκων χωρίς αλλαγή, τότε θα πρέπει να γίνει σχετική αναφορά. Αντίστροφα το σχετικό κείμενο, εικόνα, πίνακας, φυσικά, θα αναφέρονται στη βιβλιογραφία.

Οι αναφορές δεν θα πρέπει -όπως συνηθίζεται σε ορισμένες ανθρωπιστικές επιστήμες- να καταχωρίζονται σε υποσημειώσεις ή τους τίτλους των κεφαλαίων. Όταν περιγράφετε αποτελέσματα των σχετικών εργασιών, τότε πρέπει να αναφέρετε την αντίστοιχη πηγή στο κυρίως κείμενο της εργασίας σας.

\section{Υποσημειώσεις}
Χρησιμοποιήστε υποσημειώσεις στο λόγο σας με τη μεγαλύτερη δυνατή φειδώ και κρατήστε τις σύντομες.
Συχνά εμφανιζόμενες, μακροσκελής εξηγήσεις στις υποσημειώσεις διαταράσσουν τη ροή του κειμένου και το ύφος της γραφής. Το κείμενο σας πρέπει να γίνεται κατανοητό ακόμα και σε περίπτωση που δεν διαβάσετε όλες τις υποσημειώσεις.

